% GENERATED FILE. Edit canonical lessons, then run npm run generate:books.
% canonical-source-hash: fnv1a64:10be735663092866
% canonical-lessons: TA-C147-ilaiya, TA-C147-vayatana, TA-C147-putticali, TA-C147-comperi, TA-C147-tairiyamana

\chapter{Young, Old, Clever, Lazy, Brave}
\label{ch:ta-young-old-clever-lazy-brave}
\hlchaptermodality{147}

\begin{chapteropening}
\textbf{By the end of this chapter:} \emph{I can say young, old, clever, lazy and brave.}

Complete the last lesson of chapter 147: i can say young, old, clever, lazy and brave.
\end{chapteropening}

\section[\texorpdfstring{\ta{இளைய} (iḷaiya)}{இளைய (iḷaiya)}]{\ta{இளைய} (iḷaiya) — young, younger}

\label{lesson:TA-C147-ilaiya}



\begin{quote}
Before the new one: say the Tamil for rich, then the Tamil for poor.

An adjective goes before the noun and does not change: \textbf{\ta{பெரிய} \ta{வீடு}}, \textbf{\ta{சிறிய} \ta{வீடு}}.
\end{quote}

\subsection*{You'll want to know: \ta{இளைய}}
\textbf{\ta{இளைய}} — \emph{iḷaiya} — \textquotedblleft{}young, younger\textquotedblright{}.

\textbf{In the family, and next door}

\noindent
\begin{tabularx}{\linewidth}{@{}>{\raggedright\arraybackslash}X>{\raggedright\arraybackslash}X@{}}
\toprule
\textbf{} & \textbf{word} \\
\midrule
Hindi (neighbour) & \dv{जवान} (\emph{javān}) \\
English & young, younger \\
\bottomrule
\end{tabularx}

\begin{grammarlens}[title={What you've built}]
One more word for this chapter.
\end{grammarlens}

\subsection*{Your turn}
\begin{itemize}
\raggedright
  \item \emph{Say it:} \emph{iḷaiya}
  \item \emph{Say it:} \emph{iḷaiya}, once more
  \item \emph{Say it:} say \emph{iḷaiya}
  \item \emph{Recall:} say the Tamil for rich, then the Tamil for poor, then say \emph{iḷaiya} again
\end{itemize}

\begin{tcolorbox}[breakable,colback=teal!4,colframe=teal!35!black,title={Before you move on}]
What does \ta{இளைய} mean? (\textquotedblleft{}Young, younger\textquotedblright{}.) Say it once more.
\end{tcolorbox}
\section[\texorpdfstring{\ta{வயதான} (vayatāṉa)}{வயதான (vayatāṉa)}]{\ta{வயதான} (vayatāṉa) — old (of people)}

\label{lesson:TA-C147-vayatana}



\begin{quote}
Before the new one: say the Tamil for poor, then the Tamil for young.
\end{quote}

\subsection*{You'll want to know: \ta{வயதான}}
\textbf{\ta{வயதான}} — \emph{vayatāṉa} — \textquotedblleft{}old (of people)\textquotedblright{}.

\textbf{In the family, and next door}

\noindent
\begin{tabularx}{\linewidth}{@{}>{\raggedright\arraybackslash}X>{\raggedright\arraybackslash}X@{}}
\toprule
\textbf{} & \textbf{word} \\
\midrule
Kannada & \ka{ಹಳೆಯ} (\emph{haḷeya}) \\
Malayalam & \ml{പഴയ} (\emph{paḻaya}) \\
English & old \\
\bottomrule
\end{tabularx}

\begin{grammarlens}[title={What you've built}]
One more word for this chapter.
\end{grammarlens}

\subsection*{Your turn}
\begin{itemize}
\raggedright
  \item \emph{Say it:} \emph{vayatāṉa}
  \item \emph{Say it:} \emph{vayatāṉa}, once more
  \item \emph{Say it:} say \emph{vayatāṉa}
  \item \emph{Recall:} say the Tamil for poor, then the Tamil for young, then say \emph{vayatāṉa} again
\end{itemize}

\begin{tcolorbox}[breakable,colback=teal!4,colframe=teal!35!black,title={Before you move on}]
What does \ta{வயதான} mean? (\textquotedblleft{}Old (of people)\textquotedblright{}.) Say it once more.
\end{tcolorbox}
\section[\texorpdfstring{\ta{புத்திசாலி} (putticāli)}{புத்திசாலி (putticāli)}]{\ta{புத்திசாலி} (putticāli) — clever; a clever person}

\label{lesson:TA-C147-putticali}



\begin{quote}
Before the new one: say the Tamil for young, then the Tamil for old.
\end{quote}

\subsection*{You'll want to know: \ta{புத்திசாலி}}
\textbf{\ta{புத்திசாலி}} — \emph{putticāli} — \textquotedblleft{}clever; a clever person\textquotedblright{}.

\textbf{In the family, and next door}

\noindent
\begin{tabularx}{\linewidth}{@{}>{\raggedright\arraybackslash}X>{\raggedright\arraybackslash}X@{}}
\toprule
\textbf{} & \textbf{word} \\
\midrule
Kannada & \ka{ಜಾಣ} (\emph{jāṇa}) \\
Telugu & \te{తెలివైన} (\emph{telivaina}) \\
Malayalam & \ml{ബുദ്ധിയുള്ള} (\emph{buddhiyuḷḷa}) \\
Hindi (neighbour) & \dv{होशियार} (\emph{hośiyār}) \\
English & clever \\
\bottomrule
\end{tabularx}

\begin{grammarlens}[title={What you've built}]
One more word for this chapter.
\end{grammarlens}

\subsection*{Your turn}
\begin{itemize}
\raggedright
  \item \emph{Say it:} \emph{putticāli}
  \item \emph{Say it:} \emph{putticāli}, once more
  \item \emph{Say it:} say \emph{putticāli}
  \item \emph{Recall:} say the Tamil for young, then the Tamil for old, then say \emph{putticāli} again
\end{itemize}

\begin{tcolorbox}[breakable,colback=teal!4,colframe=teal!35!black,title={Before you move on}]
What does \ta{புத்திசாலி} mean? (\textquotedblleft{}Clever; a clever person\textquotedblright{}.) Say it once more.
\end{tcolorbox}
\section[\texorpdfstring{\ta{சோம்பேறி} (cōmpēṟi)}{சோம்பேறி (cōmpēṟi)}]{\ta{சோம்பேறி} (cōmpēṟi) — lazy; a lazy person}

\label{lesson:TA-C147-comperi}



\begin{quote}
Before the new one: say the Tamil for old, then the Tamil for clever; a clever person.
\end{quote}

\subsection*{You'll want to know: \ta{சோம்பேறி}}
\textbf{\ta{சோம்பேறி}} — \emph{cōmpēṟi} — \textquotedblleft{}lazy; a lazy person\textquotedblright{}.

\textbf{In the family, and next door}

\noindent
\begin{tabularx}{\linewidth}{@{}>{\raggedright\arraybackslash}X>{\raggedright\arraybackslash}X>{\raggedright\arraybackslash}X@{}}
\toprule
\textbf{} & \textbf{word} & \textbf{same root} \\
\midrule
Kannada & \ka{ಸೋಮಾರಿ} (\emph{sōmāri}) & yes \\
Telugu & \te{సోమరి} (\emph{sōmari}) & yes \\
Malayalam & \ml{മടിയുള്ള} (\emph{maṭiyuḷḷa}) &  \\
English & lazy &  \\
\bottomrule
\end{tabularx}

\begin{grammarlens}[title={What you've built}]
One more word for this chapter.
\end{grammarlens}

\subsection*{Your turn}
\begin{itemize}
\raggedright
  \item \emph{Say it:} \emph{cōmpēṟi}
  \item \emph{Say it:} \emph{cōmpēṟi}, once more
  \item \emph{Say it:} say \emph{cōmpēṟi}
  \item \emph{Recall:} say the Tamil for old, then the Tamil for clever; a clever person, then say \emph{cōmpēṟi} again
\end{itemize}

\begin{tcolorbox}[breakable,colback=teal!4,colframe=teal!35!black,title={Before you move on}]
What does \ta{சோம்பேறி} mean? (\textquotedblleft{}Lazy; a lazy person\textquotedblright{}.) Say it once more.
\end{tcolorbox}
\section[\texorpdfstring{\ta{தைரியமான} (tairiyamāṉa)}{தைரியமான (tairiyamāṉa)}]{\ta{தைரியமான} (tairiyamāṉa) — brave}

\label{lesson:TA-C147-tairiyamana}



\begin{quote}
Before the new one: say the Tamil for clever; a clever person, then the Tamil for lazy; a lazy person.
\end{quote}

\subsection*{You'll want to know: \ta{தைரியமான}}
\textbf{\ta{தைரியமான}} — \emph{tairiyamāṉa} — \textquotedblleft{}brave\textquotedblright{}.

\textbf{In the family, and next door}

\noindent
\begin{tabularx}{\linewidth}{@{}>{\raggedright\arraybackslash}X>{\raggedright\arraybackslash}X>{\raggedright\arraybackslash}X@{}}
\toprule
\textbf{} & \textbf{word} & \textbf{same root} \\
\midrule
Kannada & \ka{ಧೈರ್ಯವಂತ} (\emph{dhairyavanta}) & yes \\
Malayalam & \ml{ധൈര്യമുള്ള} (\emph{dhairyamuḷḷa}) & yes \\
English & brave &  \\
\bottomrule
\end{tabularx}

\begin{grammarlens}[title={What you've built}]
One more word for this chapter.
\end{grammarlens}

\subsection*{Your turn}
\begin{itemize}
\raggedright
  \item \emph{Say it:} \emph{tairiyamāṉa}
  \item \emph{Say it:} \emph{tairiyamāṉa}, once more
  \item \emph{Say it:} say \emph{tairiyamāṉa}
  \item \emph{Recall:} say the Tamil for clever; a clever person, then the Tamil for lazy; a lazy person, then say \emph{tairiyamāṉa} again
\end{itemize}

\begin{tcolorbox}[breakable,colback=teal!4,colframe=teal!35!black,title={Before you move on}]
What does \ta{தைரியமான} mean? (\textquotedblleft{}Brave\textquotedblright{}.) Say it once more.
\end{tcolorbox}
